\section{Algorithm Design}

Pseudo-code in this section follows the actual project scripts. It is included for algorithm design explanation and is excluded from the report word count.

\subsection{Naive Bayes}

The Naive Bayes model uses TF-IDF unigram and bigram features with \texttt{MultinomialNB}. The vectorizer is fitted only on the training split, and the trained vectorizer is then applied to the test split.

\begin{algorithm}[H]
\caption{TF-IDF + Multinomial Naive Bayes}
\label{alg:naive_bayes}
\begin{algorithmic}[1]
\Require Training texts $X_{\text{train}}$, training labels $y_{\text{train}}$, test texts $X_{\text{test}}$
\Ensure Predicted labels $\hat{y}_{\text{test}}$ and evaluation metrics
\State Fit TF-IDF unigram/bigram vectorizer on $X_{\text{train}}$ only.
\State Transform $X_{\text{train}}$ into training vectors $V_{\text{train}}$.
\State Train \texttt{MultinomialNB} on $(V_{\text{train}}, y_{\text{train}})$.
\State Transform $X_{\text{test}}$ using the fitted vectorizer to obtain $V_{\text{test}}$.
\State Predict test labels $\hat{y}_{\text{test}}$ with the trained classifier.
\State Compute precision, recall, F1-score, and accuracy.
\end{algorithmic}
\end{algorithm}

The feature representation is sparse lexical evidence. The classifier is efficient and gives a clear baseline, but its conditional independence assumption limits its ability to model word interactions.

\subsection{Support Vector Machine}

The SVM model uses TF-IDF unigram and bigram features with \texttt{LinearSVC}. A linear classifier is suitable for high-dimensional sparse text features.

\begin{algorithm}[H]
\caption{TF-IDF + Linear Support Vector Machine}
\label{alg:svm}
\begin{algorithmic}[1]
\Require Training texts $X_{\text{train}}$, training labels $y_{\text{train}}$, test texts $X_{\text{test}}$
\Ensure Predicted labels $\hat{y}_{\text{test}}$ and evaluation metrics
\State Fit TF-IDF unigram/bigram vectorizer on $X_{\text{train}}$ only.
\State Transform $X_{\text{train}}$ into training vectors $V_{\text{train}}$.
\State Train \texttt{LinearSVC} on $(V_{\text{train}}, y_{\text{train}})$.
\State Transform $X_{\text{test}}$ using the fitted vectorizer to obtain $V_{\text{test}}$.
\State Predict test labels $\hat{y}_{\text{test}}$ with the trained SVM.
\State Compute precision, recall, F1-score, and accuracy.
\end{algorithmic}
\end{algorithm}

The main component is the maximum-margin linear decision boundary. Compared with Naive Bayes, it does not model class-conditional word probabilities, but it often works well with sparse text features.

\subsection{Word2Vec-based Classifier}

The Word2Vec pipeline trains embeddings only on the training split. Each document is represented by the average of its in-vocabulary token vectors, and Logistic Regression performs final classification.

\begin{algorithm}[H]
\caption{Self-trained Word2Vec + Average Pooling + Logistic Regression}
\label{alg:word2vec}
\begin{algorithmic}[1]
\Require Training texts $X_{\text{train}}$, training labels $y_{\text{train}}$, test texts $X_{\text{test}}$
\Ensure Predicted labels $\hat{y}_{\text{test}}$ and evaluation metrics
\State Tokenize training texts into token lists $T_{\text{train}}$.
\State Train Word2Vec only on $T_{\text{train}}$.
\For{each document $d$ in training and test texts}
    \State Collect Word2Vec vectors for in-vocabulary tokens in $d$.
    \If{no token in $d$ is in vocabulary}
        \State Use a zero vector as the document vector.
    \Else
        \State Average in-vocabulary word vectors into one document vector.
    \EndIf
\EndFor
\State Train Logistic Regression on training document vectors and $y_{\text{train}}$.
\State Predict test labels $\hat{y}_{\text{test}}$ from test document vectors.
\State Compute precision, recall, F1-score, and accuracy.
\end{algorithmic}
\end{algorithm}

The implementation uses \texttt{vector\_size=100}, \texttt{window=5}, \texttt{min\_count=2}, \texttt{sg=1}, and Logistic Regression with \texttt{max\_iter=1000}. Its limitation is that average pooling removes word order and much of the local context.

\subsection{BERT-based Classifier}

The BERT-based model fine-tunes \texttt{distilbert-base-uncased} for sequence classification. Labels are encoded with \texttt{LabelEncoder}; predictions are decoded back to the original label strings before saving.

\begin{algorithm}[H]
\caption{DistilBERT Fine-tuning for Document Topic Classification}
\label{alg:bert}
\begin{algorithmic}[1]
\Require Training texts $X_{\text{train}}$, training labels $y_{\text{train}}$, test texts $X_{\text{test}}$
\Ensure Predicted labels $\hat{y}_{\text{test}}$ and evaluation metrics
\State Encode string labels as integer class IDs.
\State Load \texttt{distilbert-base-uncased} tokenizer and sequence classification model.
\State Tokenize input text with maximum length 128, truncation, and padding.
\State Fine-tune the DistilBERT classifier on the training split.
\State Evaluate the fine-tuned model on the test split.
\State Decode predicted class IDs back to original label strings.
\State Compute precision, recall, F1-score, and accuracy.
\end{algorithmic}
\end{algorithm}

DistilBERT provides contextual token representations, which helps with ambiguous or informal documents. The cost is higher training and inference time than the traditional models.
